% ---------------------------------------------------
% Energia Cinética por Elo
% ---------------------------------------------------
\section{Energia cinética (Em desenvolvimento)}
Um corpo rígido possui energia cinética pois leva em conta sua translação e sua rotação. Cada elo do manipulador robótico possui sua própria característica, como massa, centro de massa e inércia, portanto a sua energia cinética é calculada de forma individual.

\begin{equation}
T_i =
\frac{1}{2} m_i \, ({}^{\,T} \vec{v}_{C_i}) \vec{v}_{C_i}
+ \frac{1}{2} ({}^{\,T} \vec{\omega}_i) \, ({}^{(0)} \mathbf{I}_{C_i}) \, \vec{\omega}_i
\end{equation}

Onde:
\begin{itemize}
    \item $m_i$ é a massa do elo $i$;
    \item $\vec{v}_{C_i}$ é a velocidade do centro de massa do elo $i$ (em relação à base);
    \item $\vec{\omega}_i$ é a velocidade angular do elo $i$;
    \item ${}^{(0)} \mathbf{I}_{C_i}$ é o tensor de inércia do elo $i$ calculado no centro de massa, expresso no frame da base, $frame \{0\}$.
\end{itemize}

% ---------------------------------------------------
% Energia Cinética da Base
% ---------------------------------------------------
\subsection{Energia cinética da base (Em desenvolvimento)}
A base também pode contribuir com energia cinética, especialmente se houver rotação, caso haja translação da base, essa pode ser tratada separadamente como parte do subsistema orbital.

\begin{equation}
T_b = \frac{1}{2} \, \vec{\omega}_b^{\,T} \, \mathbf{I}_b \, \vec{\omega}_b
\end{equation}

Onde:
\begin{itemize}
    \item $\vec{\omega}_b$ é a velocidade angular da base;
    \item $\mathbf{I}_b$ é o tensor de inércia da base.
\end{itemize}

% ---------------------------------------------------
% Energia Cinética Total
% ---------------------------------------------------
A energia cinética total do manipulador robótico pode ser obtida ao se contabilizar todos os elos e a sua base:

\begin{equation}
T = \sum_{i=1}^{N_{\text{elos}}} T_i + T_b
\end{equation}

% ---------------------------------------------------
% Energia Potencial
% ---------------------------------------------------
Por conta da espaçonave do tipo manipulador robótico estar em ambiente de microgravidade e por não haver outro tipo de força de perturbação ambiental, é possível considerar:
\begin{equation}
V \simeq 0
\end{equation}

% ---------------------------------------------------
% Lagrangiana
% ---------------------------------------------------
Com o objetivo de encontrar o comportamento dinâmico de um sistema mecânico, utiliza-se a Lagrangiana. Ela é utilizada para derivar as equações do movimento do sistema, removendo a necessidade de analisar diretamente todas as forças que agem sobre ele, como por exemplo, a força normal, força de atrito etc.

Para encontrar a Lagrangiana do sistema, realiza-se a diferença entre a energia cinética e a energia potencial:
\begin{equation}
\mathcal{L}(q, \dot{q}) = T(q, \dot{q}) - V(q)
\end{equation}

% Ao desprezar a energia potencial por estar em ambiente de microgravidade, a Lagrangiana se reduz à própria energia cinética:
% \begin{equation}
%     \label{eq:Lagrangiana_apenasCinetica}
% \mathcal{L}(q, \dot{q}) \approx T(q, \dot{q})
% \end{equation}


